\section{Implementierung}\label{sec:implementierung}

\subsection*{Gesamtaufbau}

Die Anwendung ist eine \texttt{Browser.element}-Anwendung nach der Elm-Architektur
\cite{ElmArchitecture}: ein zentrales \texttt{Model}, ein \texttt{Msg}-Typ für alle Ereignisse, eine
reine \texttt{update}- und eine reine \texttt{view}-Funktion. Der Code ist in sechs Module
gegliedert (Tabelle~\ref{tab:module}) und folgt einer klaren Schichtung: \texttt{Data} kennt nur
Daten; die \texttt{Views}-Module kennen Daten und ihre eigene Darstellung, aber weder das
\texttt{Model} noch den \texttt{Msg}-Typ; \texttt{Main} kennt alles und verbindet es.

\begin{table}[H]
\centering\small
\caption{Modulstruktur der Anwendung (Stand Commit \texttt{\repocommit})}\label{tab:module}
\begin{tabular}{@{}llr@{}}
\toprule
Modul & Verantwortung & Zeilen \\
\midrule
\texttt{Main.elm} & Model, Msg, update, Komposition der Ansichten & 324 \\
\texttt{Data.elm} & Typen, JSON-Decoder, Filter- und Suchfunktionen, Formatierung & 342 \\
\texttt{Views/TimeSeries.elm} & Visualisierung Eins (Monatsübersicht) & 424 \\
\texttt{Views/Heatmap.elm} & Visualisierung Zwei (Stunden-Heatmap) & 277 \\
\texttt{Views/ParallelCoordinates.elm} & Visualisierung Drei (parallele Koordinaten, Brushing) & 461 \\
\texttt{Views/Layout.elm} & Filterleiste, Kennzahlen, Detailpanel, Vergleich, Tooltip, CSS & 475 \\
\midrule
\multicolumn{2}{@{}l}{Summe Elm} & 2303 \\
\texttt{scripts/fetch\_energycharts.py} & Datenabruf und Vorverarbeitung & 279 \\
\bottomrule
\end{tabular}
\end{table}

Die Schichtung wird durch ein durchgehendes Muster erzwungen: Jedes Visualisierungsmodul definiert
einen eigenen \texttt{Config msg}-Record, der den benötigten Zustandsausschnitt und die
auszulösenden Nachrichten als \emph{Funktionen} enthält. Die Module sind dadurch über \texttt{msg}
parametrisiert und kennen den konkreten \texttt{Msg}-Typ nicht:

\begin{lstlisting}[caption={Config-Record der Heatmap: Zustandsausschnitt und Callbacks},label={lst:config}]
type alias Config msg =
    { hoveredDay : Maybe String
    , selectedDays : List String
    , onHoverDay : String -> Tooltip -> msg
    , onMove : Float -> Float -> msg
    , onLeave : msg
    , onToggleDay : String -> msg
    }
\end{lstlisting}

Der Preis sind etwas längere Aufrufstellen in \texttt{Main.view}; der Gewinn ist, dass sich eine
Ansicht ohne Kenntnis der übrigen Anwendung austauschen oder wiederverwenden lässt und dass Hover-
und Klickverhalten in allen drei Ansichten garantiert identisch sind.

\subsection*{Die wichtigsten Datenstrukturen}

\paragraph{Das Model.} Der gesamte Interaktionszustand steht in einem Record; die Typentscheidungen
bilden jeweils eine Interaktionsanforderung ab.

\begin{lstlisting}[caption={Zentrales Model mit allen Interaktionszuständen},label={lst:model}]
type alias Model =
    { dataset : Result String Dataset   -- Decoder-Fehler explizit im Typ
    , selectedMonths : Set Int          -- Menge: Mehrfachauswahl von Monaten
    , hoveredDay : Maybe String         -- fluechtiger Fokus (Hover)
    , selectedDays : List String        -- dauerhafte Auswahl, Reihenfolge zaehlt
    , tooltip : Maybe Tooltip           -- Position und Inhalt am Mauszeiger
    , brushes : Dict Int ( Float, Float )  -- Achsenindex -> Wertanteil lo..hi
    , dragging : Maybe ParallelCoordinates.Drag  -- laufender Brush-Vorgang
    }
\end{lstlisting}

\texttt{Set Int} statt \texttt{Maybe Int} macht die Mehrfachauswahl von Monaten zu einer Eigenschaft
des Typs. \texttt{List String} statt \texttt{Set String} für die Tage ist Absicht: Die Reihenfolge
bestimmt, welcher Tag als zuletzt gewählter im Detailpanel erscheint. \texttt{brushes} speichert
\emph{normierte} Anteile $0..1$ statt absoluter Werte --- dadurch bleibt ein Brush gültig, wenn sich
durch einen Monatswechsel die Achsenskalierung ändert, und die Auswertung ist von der Dimension
unabhängig. \texttt{dragging} trennt den laufenden vom abgeschlossenen Brush, sodass während des
Ziehens live gefiltert werden kann, ohne den bestätigten Zustand zu überschreiben. Dass
\texttt{dataset} ein \texttt{Result} ist, zwingt \texttt{view}, den Fehlerfall zu behandeln --- ein
nicht dekodierbarer Datensatz erzeugt eine Meldung, keinen leeren Bildschirm.

\paragraph{Die Datentypen.} \texttt{Hourly} (14 Felder) und \texttt{Daily} (13 Felder) bilden die
beiden Aggregationsstufen ab. Diese Trennung ist die wichtigste Entwurfsentscheidung des
Datenmodells: Weil die Tagesaggregate vorliegen, muss keine Ansicht bei jeder Interaktion
gruppieren. \texttt{priceEurMwh : Maybe Float} macht die mögliche Abwesenheit eines Preises im Typ
sichtbar, sodass jede Ansicht sie behandeln muss.

\paragraph{Ein Detail des Decoders.} \texttt{Json.Decode} bietet Kombinatoren nur bis
\texttt{map8}; beide Records haben mehr Felder. Statt eine Zusatzbibliothek einzubinden, ist
\texttt{andMap} in drei Zeilen selbst definiert und der Decoder als Kette geschrieben:

\begin{lstlisting}[caption={andMap erweitert map8 auf beliebig viele Felder},label={lst:andmap}]
andMap : Decoder value -> Decoder (value -> result) -> Decoder result
andMap valueDecoder functionDecoder =
    Decode.map2 (\fn value -> fn value) functionDecoder valueDecoder
\end{lstlisting}

\subsection*{Funktion der drei Visualisierungen im Code}

Alle drei Ansichten erzeugen SVG mit \texttt{elm-community/typed-svg}. Der wesentliche Vorteil
gegenüber handgeschriebenem SVG ist, dass Attribute typisiert sind: Ein Längenwert ohne Einheit oder
ein falsch geschriebenes Attribut ist ein Übersetzungsfehler und kein stilles Nichtzeichnen.

\emph{Zeitreihe.} \texttt{monthly} faltet die Tagesliste zu zwölf \texttt{MonthSummary}-Records;
Monate ohne Daten werden über \texttt{filterMap} ausgelassen statt als Null gezeichnet. Die
Balkenhöhe skaliert linear auf $0..130\,\%$, die Preislinie auf das beobachtete Minimum und Maximum
--- die Preisachse ist damit datenabhängig und wird bei jedem Filterwechsel neu beschriftet.

\emph{Heatmap.} Die Zeichenfläche entsteht in drei Lagen: Stundenbeschriftungen, ein vollständiges
Raster grauer \enquote{Fehlt}-Zellen und darüber die Datenzellen. Diese Reihenfolge ist der Trick,
mit dem R9 ohne Sonderfallbehandlung erfüllt wird --- wo keine Datenzelle liegt, bleibt Grau stehen.
Die Farbfunktion bildet den Anteil auf einen Rampenparameter $t$ ab:

\begin{lstlisting}[caption={Farbrampe der Heatmap: Begrenzung bei 120\,\% statt Kappen der Daten},label={lst:color}]
renewableColor : Float -> Color.Color
renewableColor value =
    let
        t = clamp 0 1 (value / 120)
        r = 0.12 + (1 - t) * 0.72
        g = 0.25 + t * 0.47
        b = 0.42 + (1 - abs (t - 0.5) * 2) * 0.2
    in
    Color.rgb r g b
\end{lstlisting}

\emph{Parallele Koordinaten.} \texttt{dimensions} erzeugt die sieben Achsen als Records aus
Beschriftung, Zugriffsfunktion und beobachtetem Wertebereich; die Achsen sind damit datengetrieben,
eine achte Dimension wäre eine Zeile. \texttt{passesBrushes} prüft einen Tag mit \texttt{List.all}
gegen alle aktiven Brushes, was die UND-Verknüpfung direkt ausdrückt.

\subsection*{Funktion der Interaktion im Code}

Alle Ereignisse laufen über einen \texttt{Msg}-Typ mit zwölf Konstruktoren und werden in einer
einzigen \texttt{update}-Funktion behandelt. Die Kopplung entsteht dabei nicht durch Nachrichten
\emph{zwischen} Ansichten, sondern dadurch, dass alle aus demselben Zustand lesen --- das ist der
architektonische Kern der Anwendung.

Die aufwendigste Interaktion ist das Brushing, weil es Mauskoordinaten in Datenwerte übersetzen
muss. Der übliche Weg über \texttt{getBoundingClientRect} ist in Elm nur über Ports oder
\texttt{Browser.Dom}-Tasks erreichbar. Die Anwendung vermeidet das mit einer geometrischen
Entscheidung: Das SVG der parallelen Koordinaten wird in \emph{fester} Pixelgröße gerendert
($1280 \times 450$), sodass \texttt{offsetY} eines Mausereignisses direkt der SVG-$y$-Koordinate
entspricht. Damit reicht eine reine Funktion:

\begin{lstlisting}[caption={Umrechnung der Mausposition in einen normierten Achsenwert},label={lst:brush}]
yToFraction : Float -> Float
yToFraction offsetY =
    1 - clamp 0 1 ((offsetY - marginTop) / innerH)
\end{lstlisting}

Weil der Brush nicht nur die eigene Ansicht filtert, liegt seine \emph{Auswertung} nicht im
Views-Modul: \texttt{ParallelCoordinates.matchingDates} liefert die Menge der Datumswerte, die alle
gesetzten Bereiche passieren, und \texttt{Main} leitet daraus den Ausschnitt für Kennzahlen und
Heatmap ab. Der laufende Ziehvorgang geht dabei bereits ein, sodass live gefiltert wird. Das Modul
behält damit die Definition seiner Achsen --- nur dort ist bekannt, wie ein Pixelbereich in einen
Wertebereich übersetzt wird ---, gibt das Ergebnis aber als reine Funktion nach außen.

Damit \texttt{offsetY} browserübergreifend stimmt, beginnt die unsichtbare Trefferfläche jeder Achse
bei $y = 0$ und nicht erst am Achsenanfang --- andernfalls bezieht sich \texttt{offsetY} in manchen
Browsern auf die Oberkante des getroffenen Elements. Während eines Ziehvorgangs legt sich zusätzlich
ein transparentes Rechteck über die gesamte Zeichenfläche, das \texttt{mousemove}, \texttt{mouseup}
und \texttt{mouseleave} auffängt; ohne dieses Overlay bricht der Brush ab, sobald der Zeiger die
schmale Achsenspur verlässt.

\subsection*{Implementierungsaufwand}

\paragraph{Was leicht ging.} Der Aufbau der SVG-Diagramme --- Achsen, Gitter, Skalierungsfunktionen,
Pfadaufbau aus Punktlisten --- ließ sich weitgehend aus dem Code der Übungen übernehmen. Ebenso war
die Elm-Architektur durch die Übungen vertraut, sodass das Grundgerüst an einem Abend stand.
Monatsfilter und Hover-Hervorhebung sind jeweils wenige Zeilen, weil sie nur ein Feld im Model
setzen.

\paragraph{Was aufwendig war.} Erstens \emph{die Daten selbst}, nicht die Visualisierung: Dass
\texttt{load\_in\_gw} leer ist und die \enquote{\_in\_gw}-Spalten MW-Werte führen, war nur über
Plausibilitätsvergleiche zu erkennen; bis dahin waren alle Anteilswerte um Größenordnungen falsch.
Auch die Unterscheidung zwischen \texttt{where\_} und \texttt{filters} in der Schnittstelle war nur
durch systematisches Ausprobieren zu finden. Zweitens \emph{das Brushing} wegen des beschriebenen
Koordinatenproblems: Die erste Fassung arbeitete mit responsivem SVG und lieferte je nach
Fensterbreite verschobene Auswahlbereiche; die Lösung (feste Pixelgröße, Trefferfläche ab $y = 0$,
Overlay) entstand erst nach mehreren Fehlversuchen und ist im Code kommentiert, damit sie nicht
versehentlich zurückgebaut wird. Drittens \emph{die Decoder} wegen der \texttt{map8}-Grenze.
Viertens \emph{der Umgang mit Datenlücken}: Fehlende Stunden einfach wegzulassen führte dazu, dass
die Heatmap Tage stillschweigend zusammenschob und der Zeitstrahl falsch wurde --- erst die Lage aus
vorgezeichneten Grauzellen löste das strukturell.

\paragraph{Leistung.} Ungefiltert zeichnet die Anwendung rund $8\,690$ Datenzellen, $8\,784$
Hintergrundzellen und $364$ Polylinien als SVG-Elemente. Das ist im Browser flüssig, weil die
Aggregate vorberechnet sind und \texttt{update} bei jeder Interaktion nur ein Feld ändert. Für
mehrere Jahre in Stundenauflösung wäre die Grenze erreicht; dann wäre die Heatmap auf Canvas
umzustellen.

\paragraph{Fehlerbehandlung.} Die beiden wahrscheinlichsten Fehlerfälle sind getrennt behandelt:
Schlägt das \texttt{fetch()} fehl (Datei fehlt, kein Webserver), zeigt die HTML-Seite die
HTTP-Meldung; lässt sich die Datei laden, aber nicht dekodieren, zeigt Elm den Decoder-Fehler im
Klartext. Der Metadatenblock der JSON-Datei führt Land, Marktgebiet, Jahr, Quelle,
Erzeugungszeitpunkt und die Hinweise zu Normierung und Lastrekonstruktion mit.

\subsection*{Einsatz von KI-Werkzeugen}

Teile des Codes sind mit Unterstützung eines KI-Assistenten entstanden. Das Vorgehen war durchgehend
dasselbe: Wir haben Anwendungsaufgabe, gewünschte
Kodierung und Randbedingungen vorgegeben, das Ergebnis übersetzt, im Browser gegen die Daten geprüft
und gezielt nachgebessert. Selbst erstellt und nicht generiert sind die inhaltlichen Entscheidungen
dieses Berichts: die Auswahl der drei Techniken und ihre Begründung, die Achsenreihenfolge der
parallelen Koordinaten, die Behandlung der Werte über $100\,\%$ und der Datenlücken sowie die
Auswahl der Anwendungsfälle. Geprüft wurde generierter Code auf drei Wegen: Übersetzung mit
\texttt{elm make} (der Compiler schließt eine ganze Fehlerklasse aus), Vergleich der angezeigten
Werte mit unabhängig in Python berechneten Kennzahlen und manuelle Bedienung aller
Interaktionspfade.
